%%%%%%%%%%%%%%%%%%%%%%%%%%%%%%%%%%%%%%%%%
% NIH Grant Proposal for the Specific Aims and Research Plan Sections
% LaTeX Template
% Version 1.0 (21/10/13)
%
% This template has been downloaded from:
% http://www.LaTeXTemplates.com
%
% Original author:
% Erick Tatro (erickttr@gmail.com) with modifications by:
% Vel (vel@latextemplates.com)
%
% Adapted from:
% J. Hrabe (http://www.magalien.com/public/nih_grants_in_latex.html)
%
% License:
% CC BY-NC-SA 3.0 (http://creativecommons.org/licenses/by-nc-sa/3.0/)
%
%%%%%%%%%%%%%%%%%%%%%%%%%%%%%%%%%%%%%%%%%

%----------------------------------------------------------------------------------------
%	PACKAGES AND OTHER DOCUMENT CONFIGURATIONS
%----------------------------------------------------------------------------------------

\documentclass[11pt,notitlepage]{article}

% Variabili per non ripetere i contatti mille volte, MODIFICARE QUI
\newcommand{\nomeStudente}{Giovanni Bruno}
\newcommand{\cognomeStudente}{Sanna}
\newcommand{\matricolaStudente}{1029744}
\newcommand{\emailStudente}{giovannibruno.sanna@studenti.unipd.it}
\newcommand{\telStudente}{+39 0000000000}

\newcommand{\nomeTutorAziendale}{Nome}
\newcommand{\cognomeTutorAziendale}{Cognome}
\newcommand{\emailTutorAziendale}{xxtutor@azienda.it}
\newcommand{\telTutorAziendale}{+39 0000000000}

\newcommand{\ragioneSocAzienda}{Soluzioni Software S.R.L.}
\newcommand{\indirizzoAzienda}{Via Via dei Ronchi, 21, 35127 Padova (PD)}
\newcommand{\sitoAzienda}{http://www.soluzioni-sw.it/}



% A note on fonts: As of 2013, NIH allows Georgia, Arial, Helvetica, and Palatino Linotype. LaTeX doesn't have Georgia or Arial built in; you can try to come up with your own solution if you wish to use those fonts. Here, Palatino & Helvetica are available, leave the font you want to use uncommented while commenting out the other one.
%\usepackage{palatino} % Palatino font
\usepackage[utf8x]{inputenc}%codifica
\usepackage{helvet} % Helvetica font
\usepackage{multirow}
\renewcommand*\familydefault{\sfdefault} % Use the sans serif version of the font
\usepackage[T1]{fontenc}
\linespread{1.2} % A little extra line spread is better for the Palatino font
\usepackage{fancyhdr} %pacchetto per le intestazioni
\usepackage{hyperref}
\usepackage{lipsum} % Used for inserting dummy 'Lorem ipsum' text into the template
\usepackage{amsfonts, amsmath, amsthm, amssymb} % For math fonts, symbols and environments
\usepackage{graphicx} % Required for including images
%\usepackage{booktabs} % Top and bottom rules for table
%\usepackage{wrapfig} % Allows in-line images
%\usepackage[labelfont=bf]{caption} % Make figure numbering in captions bold
\usepackage[top=0.5in,bottom=0.5in,left=0.5in,right=0.5in]{geometry} % Reduce the size of the margin
\pagestyle{empty}

\hyphenation{ionto-pho-re-tic iso-tro-pic fortran} % Specifies custom hyphenation points for words or words that shouldn't be hyphenated at all

\hypersetup{
	colorlinks=true,
	linkcolor=black,
	urlcolor=blue
}

%----------------------------------------------------------------------------------------
%	Creato da Mich - Updated by Simone Pessotto 04/08/2015
%----------------------------------------------------------------------------------------

\begin{document}
	
\noindent
\parbox{0.7\columnwidth}{Università degli Studi di Padova\\
	Stage presso \ragioneSocAzienda{}\\
	\nomeStudente{} \cognomeStudente{} (\matricolaStudente{})}%matricola
\parbox{0.3\columnwidth}{
	\hfill \includegraphics[scale=0.08]{immagini/logo-unipd.png}}

\bigskip
\begin{center}
{\Huge \textbf{Resoconto Settimana 7}} \\ 
{\textbf{08/01/2018 - 12/01/2018}} \\ \bigskip
	{\Large \textit{presso \ragioneSocAzienda{}}}\\ \bigskip
	{\Large \textit{\nomeStudente{} \cognomeStudente{}}}
\end{center}

%\section*{Contatti}
%\textbf{Studente:} \nomeStudente{} \cognomeStudente{}, \href{mailto:\emailStudente{}}{\emailStudente{}}, \telStudente{} \\
%\textbf{Tutor aziendale:} \nomeTutorAziendale{} \cognomeTutorAziendale{}, \href{mailto:\emailTutorAziendale{}}{\emailTutorAziendale{}}, \telTutorAziendale{} \\
%\textbf{Azienda:} \ragioneSocAzienda{}, \indirizzoAzienda{}, \href{\sitoAzienda{}}{\sitoAzienda{}}

\section*{Scopo dello stage}

L’azienda si occupa di consulenza e sviluppo di software per la gestione aziendale, in particolare per la piccola
e media impresa. Fornisce servizi di consulenza e di riorganizzazione dei processi aziendali; inoltre si occupa
dello sviluppo di applicativi per la gestione aziendale, il CRM, la logistica e la Business Intelligence. \\ \\
Lo stage prevede l’analisi e lo sviluppo di un prototipo di ChatBot che fornisca informazioni su eventi legati alla
vita aziendale provenienti da CRM (Customer Relationship Management) o ERP (Enterprise Resource Planning). \\
Individuazione di una piattaforma di messaggistica che supporti i ChatBot e comparazione delle funzionalità
potenziali offerte; successiva realizzazione di un prototipo di ChatBot che notifichi, ad una lista di utenti definita,
eventi associati alla realtà aziendale derivanti dall’ERP o CRM. \\
Il progetto prevederà l’analisi, test e sviluppo di un prototipo di gestore di messaggi dall’azienda ai suoi dipendenti
utilizzando un ChatBot sviluppato su una delle piattaforme di messaggistica più note (WeChat, Telegram, Skype
Messenger). Lo studio individuerà una architettura che consentirà la distribuzione dei messaggi e lo sviluppo di
un prototipo che invii informazioni aziendali sia su richieste parametriche che su evento.

\bigskip
\section*{Contenuto del documento}
Nel presente documento è riportato: 
\begin{itemize}
		\item Milestone fissata per la settimana in oggetto;
		\item Riassunto del meeting settimanale avuto con il tutor aziendale;
		\item Descrizione delle attività eseguite durante la settimana;
		\item Raggiungimento risultati attesi.
	\end{itemize} 

\newpage



\noindent
\parbox{0.7\columnwidth}{Università degli Studi di Padova\\
	Stage presso \ragioneSocAzienda{}\\
	\nomeStudente{} \cognomeStudente{} (\matricolaStudente{})}%matricola
\parbox{0.3\columnwidth}{
	\hfill \includegraphics[scale=0.08]{immagini/logo-unipd.png}}

\bigskip
\section*{Milestone}
Prototipo testato e collaudato.

\bigskip
\section*{Riassunto meeting con tutor aziendale}
Di seguito sono riportati gli argomenti di discussione del meeting, avvenuto il 11/01/2018:
\begin{itemize}
		\item Presentazione del prototipo creato;
		\item Valutazione del prototipo;
		\item Risultati dei test e del collaudo effettuati;
		\item Discussione su possibili modifiche richieste dal tutor aziendale;
\end{itemize} 




\bigskip
\section*{Attività}
Di seguito sono riportate le attività svolte durante la settimana 08/01/2018 - 12/01/2018:
\begin{itemize}
		\item Implementazione dei test di unità non completati la settimana precedente;
		\item Implementazioni delle modifiche richieste dal tutor aziendale;
		\item Esecuzione dei test;
		\item Correzione degli errori emersi dai test;
		\item Collaudo del prototipo.
\end{itemize}


\newpage
\noindent
\parbox{0.7\columnwidth}{Università degli Studi di Padova\\
	Stage presso \ragioneSocAzienda{}\\
	\nomeStudente{} \cognomeStudente{} (\matricolaStudente{})}%matricola
\parbox{0.3\columnwidth}{
	\hfill \includegraphics[scale=0.08]{immagini/logo-unipd.png}}

\bigskip
\section*{Raggiungimento risultati attesi}
Questa settimana gli obiettivi prefissati sono stati completamente raggiunti; il prototipo è stato testato e collaudato. \\ Il tutor aziendale ha accettato il prototipo creato, ritenendolo molto soddisfacente.
 
 

\end{document}